\section{Realism}

很多标准 benchmark 离真实使用场景非常远。现实用户不是在做考试，而是在“问问题”。

\subsection{Quizzing vs Asking}

\begin{enumerate}
    \item \textbf{Quizzing}：用户知道答案，想测试系统。
    \item \textbf{Asking}：用户不知道答案，想借助系统得到答案。
\end{enumerate}

后者更接近真实产品使用，也更能体现模型的实际价值。

\subsection{Clio 和 MedHELM}

Clio 用模型分析真实用户数据，总结人们在问什么；MedHELM 则尝试把医疗评估从标准化考试转向更接近临床实践的任务集合。这里的核心矛盾是：越真实，越有价值，但越容易碰到隐私和合规问题。

\begin{importantbox}{真实性和隐私经常冲突}
最接近真实流量的数据，往往也是最敏感、最难共享的数据。评估如果只看标准题，会偏离使用场景；如果只看真实数据，又可能受限于隐私和合规。
\end{importantbox}

